\chapter*{Заключение}
\addcontentsline{toc}{chapter}{Заключение}

В ходе лабораторной работы выполнен сравнительный анализ классической линейной регрессии, Lasso и ElasticNet при решении задачи оценивания стоимости жилой недвижимости Египта.

\section*{Ответы на обязательные вопросы кафедры ПМ и К}

\subsection*{1. Какие основные проблемы данных были выявлены?}
При разведочном анализе исходной таблицы зафиксированы три фактора:
\begin{itemize}
    \item \textit{Пропущенные значения:} в графах этажности (\texttt{level}, 6.6\%), санузлов (\texttt{bathrooms}, 4.8\%), спален (\texttt{bedrooms}, 4.1\%) и отделки (\texttt{finishing}, 4.0\%) имелись пустые ячейки из-за необязательности этих полей при подаче объявлений;
    \item \textit{Разнородность шкал и единиц измерения:} физические размеры жилья варьируются от 40 до 650 квадратных метров, комнаты измеряются единицами, а цены исчисляются миллионами фунтов;
    \item \textit{Взаимная зависимость признаков:} число спален и санузлов закономерно растет с ростом площади, порождая частичную коллинеарность.
\end{itemize}

\subsection*{2. Какие методы предобработки были применены и почему?}
\begin{itemize}
    \item \textit{Робастная импутация пропусков:} для непрерывных величин применена медиана, нечувствительная к предельным выбросам; для категорий — наиболее частый класс (мода);
    \item \textit{Label Encoding:} категориальные атрибуты закодированы целыми числами строго по регламенту задания (п. 1.2);
    \item \textit{Стандартизация признаков (StandardScaler):} выравнивание средних к нулю и дисперсий к единице обеспечило математическую корректность $L_1$ и $L_2$ штрафов, исключив дискриминацию признаков с малым числовым диапазоном.
\end{itemize}

\subsection*{3. Как проведенная работа повлияет на дальнейшее построение моделей?}
\begin{itemize}
    \item Опытным путем доказано, что регуляризаторы со стандартными параметрами из библиотеки могут приводить к деградации качества (падение $R^2$ до 0.64 у базового ElasticNet), поэтому этап настройки сетки через \texttt{GridSearchCV} обязателен;
    \item Анализ весов выявил, что базовую стоимость формирует площадь объекта и локация, а остальные признаки служат корректирующими коэффициентами;
    \item Полученный базовый уровень качества ($R^2 \approx 0.824$) задает надежный ориентир для тестирования сложных нелинейных алгоритмов (градиентный бустинг, решающие деревья) на следующих этапах курса.
\end{itemize}

\vspace{16pt}
\noindent \textbf{Ссылка на интерактивный блокнот в Google Colab:} \\
\url{https://colab.research.google.com/github/MichaelRayven/Machine-Learning-Algorithms/blob/feat/lab4/src/lab4/lab4.ipynb}

\vspace{10pt}
\noindent \textbf{Ссылка на выбранный датасет на Kaggle:} \\
\url{https://www.kaggle.com/datasets/hassankhaled21/egyptian-real-estate-listings}
